\documentclass[11pt]{article}
\usepackage{kristiania-workbook}

\begin{document}
\begin{center}
  {\LARGE\bfseries Put your service behind HTTPS}\\[3pt]
  {\large Session 6 -- TLS, data in transit \& Wireshark}\\[5pt]
  {\small Exercise 6 \textperiodcentered\ Session 6 lab \textperiodcentered\ Estimated time: 80--110 min}
\end{center}
\vspace{2pt}\hrule\vspace{1.1em}

Write your answers in the answer document \cmd{SKY2100\_Exercise6\_Answers.docx}. Last week you ran Nextcloud over plain HTTP. Today you
\emph{see} why that is dangerous, then fix it with TLS. You will capture your own login in the clear,
turn on HTTPS, and capture again to prove the traffic is now unreadable.

\noindent\textbf{Caution:} Capture only your \emph{own} traffic in your \emph{own} VM. Sniffing other
people's traffic is both unethical and illegal. This lab is about understanding, not interception.

\wbsection{1\quad Capture the problem (plain HTTP)}
\task{Make sure Nextcloud from Exercise 5 is reachable over HTTP} (e.g.\ \cmd{http://localhost:8080}).
\task{Start a capture on the loopback (or VM) interface and log in once:}
\begin{quote}\ttfamily
sudo tcpdump -i lo -s 0 'tcp port 8080' -w http.pcap
\end{quote}
Open the capture in Wireshark (or \cmd{tcpdump -A -r http.pcap}) and find the login request.
\task{What can you read in the captured request? Quote the field names you find (do \emph{not} write a
real password here).}

\wbsection{2\quad Turn on TLS (HTTPS)}
\task{Put Nextcloud behind a reverse proxy with TLS.} The quickest option is Caddy, which obtains/serves
a certificate automatically; a self-signed certificate via Nginx is also fine. Example with Caddy:
\begin{quote}\ttfamily
sudo apt install -y caddy\\
sudo caddy reverse-proxy --from https://localhost --to localhost:8080
\end{quote}
Browse to \cmd{https://localhost} and accept the certificate (self-signed in the lab).

\wbsection{3\quad Capture again (HTTPS) and compare}
\task{Repeat the capture while logging in over HTTPS:}
\begin{quote}\ttfamily
sudo tcpdump -i lo -s 0 'tcp port 443' -w https.pcap
\end{quote}
\task{Compare the two captures. Can you find the credentials this time? What do you see instead?}

\task{In Wireshark, inspect the TLS handshake. Note the \textbf{TLS version} and that a
\textbf{certificate} is exchanged.}

\note{Capture \cmd{tcp port 8080} once more during the HTTPS login: the hop from Caddy to Nextcloud is
still plain HTTP inside your VM. TLS ends at the proxy --- keep this in mind for the Azure reflection
below.}

\aha{Put the two captures side by side. It is the same login, but in the first the field names and the
value are plain text anyone on the path could read, and in the second there is nothing but ciphertext.
Only one thing changed, TLS, and it turned readable traffic into noise. Notice what TLS does not do. It
protects the data while it travels, but the server still decrypts and reads your password at the other
end, and the capture still shows which host you talked to and when. Encryption in transit is one
specific guarantee, not a blanket of secrecy, which is why later sessions still deal with access control
and with protecting data at rest.}

\wbsection{4\quad Cloud transfer: Microsoft Learn}
\task{Work through ``Implement Azure Key Vault''} and note what a vault stores.
\par\noindent\mbox{\href{https://learn.microsoft.com/training/modules/implement-azure-key-vault/}{learn.microsoft.com/training/modules/implement-azure-key-vault}}
\task{Name the three kinds of object a Key Vault holds, and which one a TLS certificate is.}

\wbsection{5\quad Azure reflection}
\task{In Azure, TLS is often terminated at an Application Gateway using a certificate from Key Vault.
Name one advantage of terminating TLS at the gateway, and one risk it introduces for traffic
\emph{behind} the gateway.}

\wbsection{6\quad Knowledge check}
\begin{enumerate}
\item What three guarantees does TLS provide for data in transit?
\hint{Secrecy, tamper-detection, and knowing who the server really is.}
\item Difference between \textbf{symmetric} and \textbf{asymmetric} encryption (one line each).
\hint{One shared key for both directions, versus a public/private pair.}
\item Why does TLS use \emph{both} (hybrid encryption)?
\hint{Asymmetric once to agree a key, then fast symmetric for the actual data.}
\item What is a \textbf{hash} used for, and why is it not encryption?
\hint{A one-way fingerprint — it hides nothing and cannot be reversed.}
\item What does a \textbf{certificate} bind together, and who signs it?
\hint{A public key to an identity (a hostname), signed by a Certificate Authority.}
\item Explain the \textbf{chain of trust} in one or two sentences.
\hint{Server certificate to an intermediate CA to a root the browser already trusts.}
\item What attack does certificate \emph{validation} prevent, and how?
\hint{Man-in-the-middle — an attacker cannot get a valid cert for a domain it does not own.}
\item Self-signed vs CA-signed: when is each appropriate?
\hint{Self-signed for labs and internal testing; CA-signed for anything users reach.}
\item Name two \textbf{common TLS mistakes}.
\hint{Expired certs, obsolete protocols, weak ciphers, mixed content.}
\item What does \textbf{forward secrecy} protect against?
\hint{It protects yesterday's recorded traffic even if the key leaks tomorrow.}
\item Data in transit vs data at rest: give the typical control for each (CSA Domain~9).
\hint{TLS in transit; disk or storage encryption at rest.}
\item (Reflection) TLS hid your password on the wire. Name one thing TLS does \emph{not} protect.
\end{enumerate}

\signoff

\clearpage
\solheader
{\LARGE\bfseries\color{kred} Solutions}\\[2pt]
{\small Work through the lab first, then check here. Model answers are indicative — equivalent reasoning is fine.}\par\vspace{8pt}
\noindent\textbf{Note:} Model answers are indicative -- accept equivalent reasoning. Multiple-choice
answers are in \textbf{bold}.

\wbsection{Knowledge check}
\begin{enumerate}
\item \textbf{Confidentiality} (encryption), \textbf{integrity} (hashing / MAC), and \textbf{authenticity}
      (the server's certificate).
\item \textbf{Symmetric} uses one shared key for both encrypt and decrypt -- fast, but with a key-
      distribution problem. \textbf{Asymmetric} uses a public/private key pair -- it solves distribution
      but is slower.
\item TLS uses asymmetric crypto once, to authenticate the server and agree on a shared key, then fast
      symmetric encryption for the data -- the best of both (hybrid).
\item A \textbf{hash} is a one-way fingerprint used to \emph{detect change} (integrity). It is not
      encryption: it hides nothing and cannot be reversed.
\item A certificate binds a \textbf{public key} to an \textbf{identity} (e.g.\ a hostname), \textbf{signed
      by a Certificate Authority}.
\item The server certificate is signed by an intermediate CA, which is signed by a root CA trusted by
      the browser. If every signature checks out to a trusted root, the certificate is trusted (PKI).
\item \textbf{Man-in-the-middle.} An attacker cannot present a CA-signed certificate for a domain it
      does not own, so a validating client detects the imposter.
\item \textbf{Self-signed} -- labs / internal testing (encrypts, but browsers warn; no third-party
      identity). \textbf{CA-signed} -- anything public (accepted silently).
\item Any two of: expired certificates; obsolete protocols (SSLv3, TLS 1.0/1.1) or weak ciphers; mixed
      content; missing HSTS; a leaked private key.
\item \textbf{Forward secrecy} protects \emph{past} recorded sessions even if the server's long-term
      private key is later compromised (each session uses an ephemeral key).
\item In transit $\rightarrow$ \textbf{TLS}; at rest $\rightarrow$ \textbf{storage/disk encryption} with
      careful key management. (CSA Domain~9.)
\item Reflection -- TLS does not protect: metadata (that you connected and to which host); data
      \emph{after} TLS terminates; a compromised endpoint; or a malicious site that holds a valid
      certificate.
\end{enumerate}

\wbsection{Lab checkpoints}
\begin{itemize}
\item \textbf{HTTP capture:} the POST body shows the credentials in cleartext
      (\cmd{user=...\&password=...}). This is the ``aha'' moment of the lab.
\item \textbf{HTTPS capture:} only the TLS handshake (ClientHello, ServerHello, Certificate) is
      visible; the application data is encrypted (``Application Data'' records) -- no readable
      password. The TLS version should be \textbf{1.2 or 1.3}.
\item A self-signed certificate is expected to trigger a browser warning -- a good prompt to discuss
      \emph{why} validation matters.
\end{itemize}

\wbsection{References}
CSA Security Guidance v5, \textbf{Domain~9: Data Security} (verified); J\o sang -- cryptography chapter
(syllabus, page numbers not verified); Comer Ch.~17 -- Cloud Security \& Privacy (p.~233--243, context,
verified sections). MS Learn: \emph{Implement Azure Key Vault} (verified, June 2026) -- a vault stores
\textbf{keys}, \textbf{secrets} and \textbf{certificates}; a TLS certificate is a \emph{certificate}
object.

\end{document}
